% Generated draft: requires mathematical and source review.
\documentclass[11pt]{article}
\usepackage[margin=25mm]{geometry}
\usepackage{fontspec}
\setmainfont{Latin Modern Roman}
\usepackage{amsmath,amssymb,mathtools,unicode-math}
\setmathfont{Latin Modern Math}
\usepackage{enumitem,needspace,xurl,hyperref}
\setlength{\parindent}{0pt}
\setlength{\parskip}{5pt}
\setlength{\emergencystretch}{3em}
\newcommand{\N}{\mathbb{N}}
\newcommand{\Z}{\mathbb{Z}}
\newcommand{\Q}{\mathbb{Q}}
\newcommand{\R}{\mathbb{R}}
\newcommand{\C}{\mathbb{C}}
\newcommand{\F}{\mathbb{F}}
\newcommand{\A}{\mathbb{A}}
\newcommand{\PP}{\mathbb{P}}
\newcommand{\E}{\mathbb{E}}
\newcommand{\eps}{\varepsilon}
\newcommand{\dd}{\,\mathrm{d}}
\begin{document}
\section*{UnsolvedMath ID: 5100003}
\textbf{Source problem code:} catalogue entry
\subsection{Definitions and mathematical statement}
Fix an ellipse with center $O$, foci $f_1,f_2$, and a fixed confocal caustic. Let $N\geq3$ be odd and let $P=(P_1,\ldots,P_N)$ vary in a Poncelet family of $N$-periodic billiard orbits tangent to that caustic. Indices are cyclic, so $P_{N+1}=P_1$.

Let $P'$ be the outer polygon whose vertices are consecutive intersections of tangents to the original ellipse at the $P_i$, and let $P''_i$ be the caustic tangency point of the corresponding billiard side $P_iP_{i+1}$. Write $P''=(P''_1,\ldots,P''_N)$ and let $A,A',A''$ denote the respective areas of $P,P',P''$. The source also uses $\theta_i$ for the angle of $P$ at $P_i$,
\[
L=\sum_{i=1}^N\|P_{i+1}-P_i\|_2,\qquad
l_i=\|P''_i-P_i\|_2,\qquad r_i=\|P_{i+1}-P''_i\|_2.
\]
These auxiliary quantities are part of the supplied setup but do not occur in this invariant.

Invariant $k_{109}$ asserts that there is a constant $c$, depending on the fixed family but not on its member $P$, such that
\[
\frac{A(P)}{A''(P)}=\frac{A'(P)}{A(P)}=c
\]
for every member where both ratios are defined. Thus equality of the two ratios and constancy throughout the family are both requested. The supplied solved-status explanation cites the other invariants $k_{103}=A'/A$ and $k_{112}=A'A''/A^2=1$.

The record does not specify signed versus ordinary area for self-intersecting orbits, the allowed caustic type, or degeneracies with zero area or parallel consecutive tangents. The displayed ratios require $A\ne0$, $A''\ne0$, and defined outer vertices. These conditions identify necessary domains rather than an undocumented claim that every source orbit satisfies them.
\subsection{Short English statement}
For an odd-period elliptic billiard Poncelet family, the orbit-to-inner-polygon area ratio should be constant and equal to the outer-polygon-to-orbit area ratio. The supplied summary reports this as a consequence of other proven invariants.
\subsection{Sources}
\begin{itemize}
\item Local UnsolvedMath registry, permanent ID 5100003, exact source code catalogue entry, ``Elliptic-billiard invariant $k_{109}$''; proposed by Dan Reznik, Ronaldo Garcia, and Jair Koiller, 2021. Registry status solved, classification SOLVED-IN-LITERATURE.
\item Reznik, Garcia, and Koiller, \emph{Eighty New Invariants of N-Periodics in the Elliptic Billiard} (2021), table row $k_{109}$, \url{https://arxiv.org/abs/2004.12497}. Supplied source-TeX extraction accessed July 29, 2026.
\item The supplied research summary reports $k_{109}$ resolved using proven $k_{103}$ and $k_{112}$, the latter attributed to reference [3]. The record does not reproduce the bibliography entry for [3] or the precise area/caustic hypotheses of these results. Historical-source provenance says the item is presented as open and records no later authoritative resolution check; the claimed solution is therefore preserved as supplied scope-dependent status evidence.
\item Provenance flags status and rights NEEDS\_REVIEW: public access was verified, but redistribution permission was not established.
\end{itemize}
\end{document}
